\documentclass[12pt]{article}
\usepackage[utf8]{inputenc}
\usepackage[spanish]{babel}
\usepackage{tcolorbox}
\usepackage{geometry}
\geometry{margin=1.5cm}
\usepackage{xcolor}
\usepackage{titlesec}
\usepackage{multicol}

% Colores personalizados
\definecolor{CanvasBlue}{RGB}{220,230,241}
\definecolor{CanvasGray}{RGB}{240,240,240}
\definecolor{CanvasHeader}{RGB}{60,90,150}

% Estilo de títulos
\titleformat{\section}{\Large\bfseries\color{CanvasHeader}}{}{0pt}{}
\titleformat{\subsection}{\bfseries\color{CanvasHeader}}{}{0pt}{}

\begin{document}

\begin{center}
    \LARGE \textbf{Machine Learning Canvas}\\[6pt]
    \large \textbf{Proyecto: Predicción de popularidad de noticias online}
\end{center}

\vspace{0.5cm}

% Estructura del Canvas en bloques
\begin{tcolorbox}[colback=CanvasBlue!20,colframe=CanvasHeader,title=\textbf{Prediction Task}]
El objetivo es predecir la popularidad de artículos periodísticos en línea antes de su publicación, utilizando como variable objetivo el número de veces que fueron compartidos en redes sociales (\texttt{shares}). 
Se trata de un problema de \textbf{regresión supervisada}, donde se busca estimar un valor numérico continuo.
\end{tcolorbox}

\begin{tcolorbox}[colback=CanvasBlue!20,colframe=CanvasHeader,title=\textbf{Decisions}]
El modelo permitirá a los equipos editoriales y de marketing tomar decisiones basadas en datos, como:
\begin{itemize}
    \item Destacar artículos con mayor probabilidad de viralizarse.
    \item Redirigir esfuerzos publicitarios hacia contenidos más prometedores.
    \item Priorizar temas o categorías con alto nivel de interacción.
\end{itemize}
Estas decisiones contribuirán a optimizar los recursos de publicación y mejorar el impacto del contenido.
\end{tcolorbox}

\begin{tcolorbox}[colback=CanvasBlue!20,colframe=CanvasHeader,title=\textbf{Value Proposition}]
La implementación de modelos de Machine Learning ofrece a los medios digitales la posibilidad de anticipar el éxito de sus publicaciones, con beneficios como:
\begin{itemize}
    \item Asistencia automatizada en la selección y promoción de contenidos.
    \item Reducción de tiempo y costos en campañas inefectivas.
    \item Identificación de los factores que influyen en la viralidad de los artículos.
\end{itemize}
\end{tcolorbox}

\begin{tcolorbox}[colback=CanvasGray,colframe=CanvasHeader,title=\textbf{Data Sources}]
El proyecto emplea el dataset \textbf{Online News Popularity (UCI)}, que contiene 61 atributos, de los cuales 58 son predictivos, 2 no predictivos (\texttt{url}, \texttt{timedelta}) y 1 variable objetivo (\texttt{shares}). 

Incluye dimensiones como estructura del texto, análisis de sentimiento, temática, temporalidad y rendimiento en redes sociales.

Durante el preprocesamiento se eliminaron duplicados, se verificó la ausencia de valores nulos y se normalizaron variables booleanas.
\end{tcolorbox}

\columnbreak

\begin{tcolorbox}[colback=CanvasGray,colframe=CanvasHeader,title=\textbf{Features}]
El modelo considera distintas características del artículo agrupadas en cinco categorías principales:
\begin{itemize}
    \item \textbf{Texto y tono:} variables como \texttt{n\_tokens\_content} o \texttt{global\_subjectivity}, que reflejan la extensión y el sentimiento del contenido.
    \item \textbf{Estructura y medios:} número de imágenes, vídeos y enlaces, asociados al atractivo visual.
    \item \textbf{Temática:} canal de publicación (\texttt{data\_channel\_is\_tech}, \texttt{is\_world}, etc.), que influye en el interés del lector.
    \item \textbf{Temporalidad:} día de la semana y variable \texttt{is\_weekend}, útiles para detectar patrones de publicación.
    \item \textbf{Título y keywords:} polaridad y fuerza semántica del título, relacionadas con su potencial de viralidad.
\end{itemize}
\end{tcolorbox}

\begin{tcolorbox}[colback=CanvasBlue!20,colframe=CanvasHeader,title=\textbf{Building Models}]
El enfoque del proyecto es de \textbf{regresión supervisada}. Se entrenarán y compararán los siguientes modelos:
\begin{itemize}
    \item Elastic Net
    \item MLPerceptron
    \item Random Forest
    \item HistGradientBoosting
\end{itemize}
\textbf{Evaluación:} las métricas utilizadas serán MSE, MAE y R\textsuperscript{2}.
\end{tcolorbox}

\begin{multicols}{2}
\begin{tcolorbox}[colback=CanvasGray,colframe=CanvasHeader,title=\textbf{Making Predictions}]
Las predicciones se realizarán en Google Colab usando el dataset limpio. 
Se compararán los valores reales y estimados de \texttt{shares} para evaluar el rendimiento del modelo.
\end{tcolorbox}

\columnbreak

\begin{tcolorbox}[colback=CanvasGray,colframe=CanvasHeader,title=\textbf{Impact Simulation}]
Resultados esperados:
\begin{itemize}
    \item Determinar las características más influyentes.
    \item Lograr un R\textsuperscript{2} superior a 0.6 y bajo MAE.
    \item Evaluar el desempeño comparativo de los modelos.
\end{itemize}
\end{tcolorbox}
\end{multicols}

\begin{tcolorbox}[colback=CanvasBlue!20,colframe=CanvasHeader,title=\textbf{Monitoring}]
Aunque aún no se cuenta con despliegue continuo, se usará DVC y GitHub para versionar datasets y modelos, registrar métricas y asegurar reproducibilidad.
\end{tcolorbox}

\end{document}
